\scanpage{208}
\deleted{\unsure{Le|4.}} Appelons ``prégraphe'' \unsure{un|une} système

\[
\Gamma=(S,S^\vee,R\subset S\times S^\vee) \tag*{(56)}
\]

de deux ensembles $S,S^\vee$, et d'une relation entre eux -- il sera dit \uline{séparant} si les applications

\[
S\mathrel{\unsure{\rule{1.35em}{0.35pt}|\longrightarrow}}\mathscr P(S^\vee),\qquad S^\vee\mathrel{\unsure{\rule{2em}{0.35pt}|\longrightarrow}}\mathscr P(S)
\]

correspondantes sont injectives. Le symétrique du prégraphe (56) est le prégraphe

\[
\Gamma^\vee=(S^\vee,S,{}^sR\subset S^\vee\times S). \tag*{(57)}
\]

On appelle ``\uline{homothétie}'' du prégraphe toute paire de permutations

\[
\lambda=(\lambda_S,\lambda_{S^\vee})\in\unsure{\mathfrak S|\mathscr S}_S\times\mathfrak S_{S^\vee} \tag*{(58)}
\]

satisfaisant la condition

\[
\left\{\begin{aligned}
&(\lambda_S,\lambda_{S^\vee}^{-1})(R)=R\qquad\text{i.e.}\\
&\forall(s,s^\vee)\in S\times S^\vee,\ \text{on a }\bigl((s,s^\vee)\in R\iff(\lambda s,\lambda^{-1}s^\vee)\in R,
\end{aligned}\right. \tag*{(59)}
\]

\deleted{\unsure{Ce|Cela|elles}} sont donc les éléments \deleted{de $(\lambda,\lambda')$} \deleted{$\unsure{(s,t)\mapsto}$} $(u,u')\in\mathfrak S_S\times\mathfrak S_{S^\vee}$ tels que l'on ait

\[
\left\{\begin{aligned}
&(u,u'^{-1})\in\operatorname{Aut}(S,S',R)\qquad\text{i.e.}\\
&(u^{-1},u')\in\operatorname{Aut}(S,S',R)
\end{aligned}\right. \tag*{(60)}
\]

Notons que si $R$ est séparante en $S^\vee$, alors $\lambda_{S^\vee}$ est \unsure{connu|connue} quand on connaît $\lambda_S$, \unsure{celui|qui|elle|et} est défini à partir de $\lambda_S$ par

\[
\deleted{\unsure{\lambda_S(s^\vee)|\lambda_{S^\vee}(s^\vee)}}\
\left\{\begin{aligned}
&{}^sR(\lambda_{S^\vee}(s^\vee))=\lambda_S^{-1}(R(s^\vee))\qquad\text{ou aussi}\\
&{}^sR(\lambda_{S^\vee}^{-1}(s^\vee))=\lambda_S(R(s^\vee))
\end{aligned}\right. \tag*{(61)}
\]

et inversement, si $R$ est sép en $S$, $\lambda_{S^\vee}$ détermine $\lambda_S$.

% archive sheet number: 207
